\begin{definition}
    Непустое множество L называется \textit{линейным пространством}, если оно удовлетворяет таким условиям
    \begin{enumerate}
        \item Для любых $x, y \in L$ однозначно определен третий элемент $z \in L$, называемый их суммой, причем \begin{enumerate}
            \item $x + y = y + x$
            \item $x + (y + z) = (x + y) + z$
            \item $\exists 0 \in L \; \forall x \in L:\; x + 0 = x$
            \item $\forall x \in L \; \exists (-x) \in L:\; x + (-x) = 0$
        \end{enumerate}
        \item Для любого числа $\alpha$ и любого элемента $x \in L$ определен элемент $\alpha x \in L$, называемый произведением, причем
        \begin{enumerate}
            \item $\alpha (\beta x) = (\alpha \beta) x$
            \item $1 \cdot x = x$
            \item $(\alpha + \beta) x = \alpha x + \beta x$
            \item $\alpha (x + y) = \alpha x + \alpha y$
        \end{enumerate}
    \end{enumerate}
\end{definition}

\begin{definition}
    Линейные пространства $L$ и $L^*$ называются \textbf{изоморфными}, если между их элементами можно установить взаимно однозначное соответствие $f$, которое согласуется с операциями в $L$ и $L^*$:
    \[
    \forall x, y \in L,\; \forall x^*, y^* \in L^*,\; x \xleftrightarrow[f]{} x^*,  y \xleftrightarrow[f]{} y^* \implies x + y \xleftrightarrow[f]{} x^* + y^*, \alpha x  \xleftrightarrow[f]{} \alpha x^*
    \]
\end{definition}

\begin{definition}
    Евклидовы пространства $E$ и $E^*$ называются \textbf{изоморфными}, если они изоморфны как линейные пространства и при этом сохраняется скалярное произведение:
     \[
    \forall x, y \in E,\; \forall x^*, y^* \in E^*,\; x \xleftrightarrow[f]{} x^*,  y \xleftrightarrow[f]{} y^* \implies (x, y) = (x^*, y^*).
    \]
\end{definition}

\begin{note}
    Смежные опеределения сформулированы в пункте~\nameref{seventh-bilet}.
\end{note}

\begin{example}
    Определим пространство $X = \{u \in C^1[0, 1] : u(0) = 0\}$.\\
    
    Введем нормы $||u||_1 = (\int_0^1(u(x)^2 + u'(x)^2)dx)^{\frac{1}{2}}, ||u||_2 = (\int_0^1u'(x)^2dx)^{\frac{1}{2}}$. Докажем, что нормы эквивалентны. Ясно, что $||u||_2 \leqslant ||u||_1$, поэтому достаточно доказать противоположное неравенство.

    В силу условия $u(0) = 0$ справедливо равенство $$u(x) = \int_0^xu'(t)dt \quad \forall x \in [0, 1]$$

    Возведем обе части равенства в квадрат и воспользуемся неравенством Коши-Буняковского: $$u^2(x) = \left(\int_0^xu'(t)dt \right)^2 \leqslant \int_0^xdt \cdot \int_0^x u'(t)^2dt \leqslant x \int_0^1u'(t)^2dt \leqslant \int_0^1u'(t)^2dt \quad \forall x \in [0, 1]$$

    Отсюда легко вывести два следующих неравенства:
    $$\left(\int_0^1u^2(x)dx\right)^{\frac{1}{2}} \leqslant \left(\int_0^1 u'(x)^2dx\right)^{\frac{1}{2}} \quad \forall u \in X$$
    $$\underset{0 \leqslant x \leqslant 1}{max}|u(x)| \leqslant \left(\int_0^1 u'(x)^2dx\right)^{\frac{1}{2}} \quad \forall u \in X$$

    Из предпоследнего следует, что $||u||_1 \leqslant \sqrt{2}||u||_2$.
\end{example}

\begin{example}
    Теперь пример не эквивалентных норм.
    Рассмотрим множество непрерывных функций на $[a, b]$ и нормы $||x||_C = \underset{a\leqslant t\leqslant b}{max}|x(t)|, ||x||_1 = \int_a^b|x(t)|dt$. Докажем, что они не эквивалентны. На самом деле, в одну сторону нужное неравенство даже работает -- $||\cdot||_C$ сильнее, чем $||\cdot||_1$. А вот в обратную сторону не получается. 
    
    Для доказательства этого построим последовательность функций $x_n$, которые возрастают линейно от $0$ до $2n$ на отрезке $[a, a+\frac{1}{2n}]$, убывают линейно от $2n$ до $0$ на отрезке $[a+\frac{1}{2n}, a+\frac{1}{n}]$ и равны нулю в остальных точках отрезка. Тогда $||x||_1 = 1, ||x||_C = n \to \infty$. Поэтому не существует такой константы $K$, для которой было бы верно 
    $$||x_n||_C \leqslant K \cdot \int_a^b|x_n(t)|dt$$
\end{example}

\begin{proposition}\label{normequiv}
	Пусть $E$ "--- линейное нормированное пространство, и $\dim E < +\infty$. Тогда любые две нормы на $E$ эквивалентны.
\end{proposition}

\begin{proof}(Доказательство для случая, когда {$\Kk = \R$})
	Поскольку $E$ конечномерно, то в нем можно выбрать максимальную по включению линейно независимую систему $e = \{e_1, \dotsc, e_n\} \subset E$, тогда $e$ будет являться базисом в $E$. Для произвольного элемента $x \in E$, имеющего в базисе $e$ координатный столбец $\alpha \in \R^n$, зададим его \textit{евклидову норму} следующим образом:
	\[\|x\|_{euc} := \sqrt{\sum_{k = 1}^n \alpha_k^2}\]

	\noindentЗафиксируем произвольную норму $\|\cdot\|$ и докажем, что она эквивалентна норме $\|\cdot\|_{euc}$
	\begin{enumerate}
		\item Покажем, что $\|\cdot\| < C \|\cdot\|_{euc}$ для некоторого $C > 0$. Для произвольного $x \in E$, имеющего в базисе $e$ координатный столбец $\alpha \in \R^n$, выполнено следующее:
		\[\|x\| = \left\|\sum_{k = 1}^n\alpha_k e_k \right\| \leqslant \max_{1 \leqslant k \leqslant n}\|e_k\|\left(\sum_{k = 1}^n|\alpha_k|\right)\]
  
		Поскольку для любого $k \in \{1, \dotsc, n\}$ выполнено $|\alpha_k| < \|x\|_{euc}$, то достаточно взять число $C := n(\max_{1 \leqslant k \leqslant n}\|e_k\|)$.

		\item Покажем теперь, что $\|\cdot\|_{euc} < \widetilde C \|\cdot\|$ для некоторого $\widetilde C > 0$. Предположим противное, тогда для любого $n \in \N$ существует $x_n \in E$ такое, что $\|x_n\|_{euc} > n\|x_n\|$. Можно без ограничения общности считать, что $\|x_n\|_{euc} = 1$ для любого $n \in \N$, откуда $\|x_n\| < \frac 1n$.
		
		Поскольку последовательность $\{x_n\}$ содержится в единичной сфере $S_{euc}(0, 1)$ и относительно евклидовой нормы сфера компактна, можно выделить из $\{x_n\}$ подпоследовательность $\{x_{n_k}\}$, сходящуюся относительно евклидовой нормы. Тогда $\|x_{n_k} - x\|_{euc} \to 0$ для некоторого $x \in S_{euc}(0, 1)$, тогда, в силу пункта $(1)$, $\|x_{n_k} - x\| \to 0$. Но по построению $\|x_{n_k}\| < \frac 1{n_k}\to 0$, поэтому $x = 0$ --- противоречие с тем, что $x \in S_{euc}(0, 1)$.\qedhere
	\end{enumerate}
\end{proof}

\begin{corollary}\label{closeness-subspace}
	Пусть $E$ "--- линейное нормированное пространство, $x_1, \dotsc, x_n \in E$. Тогда линейная оболочка $L := [x_1, \dotsc, x_n]$ образует подпространство в $E$.
\end{corollary}

\begin{proof}
	Заметим, что $\dim L < +\infty$, и по утверждению \ref{normequiv} сужение нормы из $E$ на $L$ эквивалентно евклидовой норме. Относительно евклидовой нормы конечномерное пространство полно, поэтому и $L$ полно относительно нормы из $E$. Следовательно, $L$ замкнуто как подмножество в $E$.
\end{proof}

\begin{definition}
	\textit{Линейным топологическим пространством} называется топологическое пространство $X$ с определенными на нем операциями сложения и умножения на числа из поля $\K$, непрерывными на $X$.
\end{definition}

\begin{example}
	Любое линейное нормированное пространство $E$ является линейным топологическим, как и пространство основных функций $D$, которое при этом даже не является метризуемым.
\end{example}
